\begin{itemize}
% Colocación y evaluación de alarmas según sensibilidad
\item Coloque su alarma en posición de alta sensibilidad, su detector en la escala más baja e inicie la búsqueda a la brevedad posible.

% Acciones de evacuación del área
\item Evacue el área de trabajo.

% Notificación a entidades relevantes
\item Notifique a las oficinas generales y a la C.N.S.N.S.

% Revisión de placas para verificar si la fuente se perdió
\item Revise las placas tomadas para determinar si la fuente se perdió desde el inicio de la jornada de trabajo, esto le ayudará a precisar el lugar más probable de su pérdida.

% Procedimientos en caso de pérdida de una fuente radiactiva
\item Tan pronto la localice proceda a su recuperación y evalúe las posibles personas que pudieron estar expuestas durante el tiempo que permaneció la fuente perdida.

% Preparación de informes detallados sobre exposición 
\item Obtenga un informe detallado de los tiempos y tipos de maniobras que ejecutó cada una de las personas con sospecha de sobreexposición y actúe conforme al punto 7.10.
\end{itemize}
